% =====================================================================
%  CAPÍTULO 1 — INTRODUCCIÓN
% =====================================================================
\chapter{Introducción}
\label{cap:introduccion}

\section{Contexto y justificación}
\label{sec:contexto}

El 5 de enero de 2023 entró en vigor la Directiva (UE) 2022/2464, conocida como
\textit{Corporate Sustainability Reporting Directive} (CSRD), y con ella se abrió la que debía ser una de
las transformaciones más profundas que ha vivido la información corporativa europea en las
últimas décadas. Sin embargo, su efecto no se ha percibido hasta hace muy poco:
el ejercicio 2024 ha sido el primero en el que las grandes empresas de interés público
europeas han estado obligadas a reportar su desempeño en sostenibilidad conforme a un marco
único, los \textit{European Sustainability Reporting Standards} (ESRS), cuyos primeros
informes se han publicado a lo largo de 2025. Por primera vez, miles de compañías han tenido
que aplicar el principio de doble materialidad, cuantificar sus emisiones a lo largo de toda
la cadena de valor y someter parte de esa información a verificación externa. El reporting de
sostenibilidad ha dejado de ser un ejercicio mayoritariamente voluntario y promocional para
convertirse en una obligación regulada, comparable y auditable \parencite{csrd2022}.

Este giro normativo se produce en un contexto en el que la información no financiera ha
ganado un peso enorme. Inversores institucionales, supervisores y consumidores
demandan saber no solo cuánto gana una empresa, sino cómo lo gana: qué impacto ambiental y
social genera, qué riesgos de transición asume y qué compromisos adquiere de cara al futuro.
La sostenibilidad se ha incorporado al núcleo de la estrategia y de la comunicación
corporativas, y los informes anuales dedican un espacio creciente a explicarla. Ahora bien,
esa misma centralidad ha hecho aflorar un problema bien conocido en la literatura: la
distancia entre lo que las empresas \textit{dicen} y lo que realmente \textit{hacen}. El
fenómeno del \textit{greenwashing} —la construcción de una imagen de sostenibilidad más
favorable que el desempeño efectivo— y su versión más sutil, el \textit{cheap talk}
—el compromiso retórico que no se traduce en acciones concretas ni en datos verificables—,
se han convertido en una de las grandes preocupaciones regulatorias y académicas del momento
\parencite{bingler2022}.

Aquí reside el problema práctico que motiva este trabajo. Buena parte de los análisis de
sostenibilidad descansan sobre las calificaciones o \textit{ratings} ESG que elaboran
proveedores externos como MSCI, Sustainalytics o Refinitiv. Estos indicadores son cómodos,
pero presentan dos limitaciones serias: por un lado, son poco transparentes y muestran
divergencias notables entre proveedores para una misma empresa; por otro, y de manera más
fundamental para esta investigación, introducen una circularidad metodológica si lo que se
pretende es precisamente evaluar la calidad de la comunicación. Si el objetivo es estudiar
\textit{cómo} comunica una empresa su sostenibilidad y detectar señales de discurso vacío,
emplear una nota externa como medida de calidad equivale a dar por respondida la pregunta
antes de formularla. Por ello, el presente trabajo renuncia deliberadamente a cualquier
\textit{rating} de terceros y se construye íntegramente sobre el análisis del \textbf{texto}
de los informes: son las propias palabras de las empresas las que se interrogan.

Para hacerlo a escala, este TFG recurre a técnicas de procesamiento del lenguaje natural
(PLN), una rama de la inteligencia artificial que permite analizar grandes volúmenes de
texto de forma sistemática y reproducible. La combinación de un marco regulatorio en plena
transición, un problema de relevancia económica como el \textit{greenwashing} y un
instrumental analítico capaz de leer miles de páginas de informes resulta especialmente
oportuna: permite observar, casi en tiempo real, cómo cambia el lenguaje corporativo de la
sostenibilidad en el tránsito de un régimen voluntario a uno obligatorio.

\bigskip
% --- NOTA: trayectoria personal inventada (placeholder) — sustituir por la real ---
La elección del tema responde también a una motivación personal. A lo largo del grado, lo
que más me ha atraído ha sido el punto en el que el análisis de datos se cruza con las
decisiones empresariales: la idea de que detrás de un texto aparentemente cualitativo, como
es un informe corporativo, se esconden patrones medibles que pueden contar una historia.
Desde que tuve mi primer contacto con la minería de datos comprendí que las herramientas
para extraer esa información ya existen y están al alcance de cualquiera con la curiosidad
suficiente para aprenderlas. Cuando se acercó el momento de plantear el TFG me pareció una
oportunidad inmejorable para llevar esa inquietud a un problema económico y financiero real,
de plena actualidad y con implicaciones tangibles para empresas, inversores y reguladores.
Aplicar la inteligencia artificial a un terreno tradicionalmente reservado al análisis manual
—la lectura de cientos de informes de sostenibilidad— era, para mí, la manera más honesta de
cerrar el grado: poniendo a trabajar juntas las dos cosas que más me interesan.

\bigskip
Este trabajo se justifica, así, por su oportunidad temporal —captura el momento
exacto de la transición normativa—, por su enfoque metodológico —análisis textual directo,
sin intermediación de notas externas— y por su contribución: ninguna investigación previa,
hasta donde alcanza la revisión realizada, ha analizado el conjunto del STOXX Europe 600 bajo
el marco completo CSRD/ESRS combinando modelado de temas, análisis de sentimiento y detección
de señales de \textit{greenwashing} sobre un corpus propio.

\section{Objetivos}
\label{sec:objetivos}

El \textbf{objetivo general} de este Trabajo Fin de Grado es analizar, mediante técnicas de
procesamiento del lenguaje natural, cómo las grandes empresas europeas del índice STOXX
Europe 600 comunican su estrategia de sostenibilidad en sus informes corporativos bajo el
marco normativo CSRD/ESRS, e identificar en ese discurso patrones comunicativos y señales
textuales de \textit{greenwashing}.

Este objetivo general se concreta en cuatro objetivos específicos, formulados como preguntas
de investigación (RQ, por sus siglas en inglés) que estructuran todo el análisis posterior:

\begin{itemize}
  \item \textbf{OE1 (RQ1) — Temas predominantes.} Identificar qué temas concentran el
  discurso de sostenibilidad de las empresas del STOXX 600, mediante modelado de temas sobre
  el texto de sus informes, y contrastar esos temas con las categorías de los estándares
  ESRS.
  \item \textbf{OE2 (RQ2) — Diferencias por sector y país.} Determinar si existen diferencias
  estadísticamente significativas en el contenido y el tono del reporting en función del
  sector de actividad y del país de origen de la empresa.
  \item \textbf{OE3 (RQ3) — Determinantes y señales de \textit{greenwashing}.} Analizar qué
  factores (sector, país, tamaño) se asocian a una mayor especificidad y cobertura
  cuantitativa de la información, y contrastar la hipótesis de que un tono más optimista se
  corresponde con un discurso menos específico —una señal característica del
  \textit{greenwashing}—.
  \item \textbf{OE4 (RQ4) — Evolución NFRD $\rightarrow$ CSRD.} Examinar cómo evoluciona el
  reporting de sostenibilidad entre el régimen de la Directiva sobre información no financiera
  (NFRD, ejercicios 2022 y 2023) y el primer ejercicio obligatorio de la CSRD (2024).
\end{itemize}

Esta precisión en la formulación de los objetivos tiene una finalidad concreta: las conclusiones del
trabajo (capítulo~\ref{cap:conclusiones}) se construyen de
forma directa sobre cada uno de ellos, de modo que el lector pueda verificar que todas las
preguntas planteadas al inicio reciben respuesta.

\section{Metodología y fuentes de información}
\label{sec:metodologia-resumen}

El trabajo adopta un enfoque cuantitativo de análisis de contenido textual asistido por
inteligencia artificial. La unidad de análisis es la subsección de sostenibilidad de cada
informe corporativo, y el diseño es un panel que combina la dimensión transversal —empresas
de distintos sectores y países— con la temporal —tres ejercicios consecutivos—. El detalle
completo se expone en el capítulo~\ref{cap:metodologia}; aquí se ofrece
únicamente una panorámica que permita situar al lector.

En cuanto a las \textbf{fuentes de información}, se distingue entre fuentes primarias y
secundarias. La fuente primaria es un corpus documental construido específicamente para esta
investigación: los informes integrados y anuales, en formato PDF, de una muestra de
\textbf{196 empresas} del STOXX Europe 600, correspondientes a los ejercicios 2022, 2023 y
2024. Tras la extracción y depuración del texto, ese corpus quedó conformado por
\textbf{586 documentos} analizables, de los que se aislaron tanto el informe de gestión como
la subsección de sostenibilidad, dando lugar a un corpus de 1.172 unidades de texto
que, una vez segmentadas, suman aproximadamente 245.400 párrafos y cerca de
540.000 frases. Como fuentes secundarias se emplearon los datos financieros de las
empresas —obtenidos de Yahoo Finance \parencite{yfinance2024}—, los textos normativos
europeos \parencite{nfrd2014,csrd2022,esrs2023,taxonomia2020,sfdr2019} y diversos diccionarios
y recursos lingüísticos especializados.

Sobre ese corpus se aplicó una batería de técnicas de PLN que van desde los diccionarios
clásicos de análisis textual financiero \parencite{loughran2011} hasta modelos de lenguaje
recientes especializados en finanzas y clima \parencite{araci2019,webersinke2022}, pasando
por dos métodos complementarios de modelado de temas \parencite{blei2003,grootendorst2022}.
A partir de sus resultados se construyó un índice sintético de \textit{greenwashing} textual
y se realizaron diversos contrastes estadísticos. Todo el proceso se organizó como un flujo
de trabajo reproducible, estructurado en siete fases sucesivas que van de la construcción de
la muestra a la redacción de este documento.

Resta una decisión metodológica que atraviesa todo el trabajo y que ya
se ha avanzado: en coherencia con su objeto de estudio, el análisis prescinde por completo de
calificaciones ESG de terceros. Toda medida de calidad, tono o especificidad procede del
propio texto de los informes, lo que preserva la independencia entre el objeto analizado
—la comunicación— y el instrumento de medida.

\section{Orden documental}
\label{sec:orden}

El presente documento se organiza en seis capítulos. Tras esta \textbf{introducción}, que
delimita el contexto, los objetivos y el enfoque general, el \textbf{capítulo~2} desarrolla el
marco del trabajo. Este se aborda desde una doble perspectiva: por un lado, el marco
normativo, que recorre la evolución de la regulación europea de la información de
sostenibilidad desde la NFRD hasta la CSRD, los ESRS, la Taxonomía de la UE y el reglamento
SFDR; por otro, el marco teórico, que recoge las aportaciones de la literatura sobre
legitimidad organizativa, gestión de impresiones y \textit{greenwashing}, así como el estado
de la cuestión en la aplicación del PLN al reporting de sostenibilidad.

El \textbf{capítulo~3} detalla la metodología: el diseño de la investigación y las preguntas
que la guían, la construcción de la muestra y del corpus, las técnicas de análisis textual
empleadas —con cita expresa de la fuente original de cada una—, la construcción del índice de
\textit{greenwashing} y el aparato estadístico utilizado para contrastar las hipótesis.

El \textbf{capítulo~4} presenta y discute los resultados, organizados en torno a las cuatro
preguntas de investigación: la radiografía del corpus y su cobertura de los estándares ESRS,
los temas predominantes, el tono y la especificidad del discurso, el índice de
\textit{greenwashing}, los factores que lo determinan y, finalmente, la evolución del
reporting entre los dos regímenes normativos. Cada hallazgo se pone en relación con el marco
teórico del capítulo~2, señalando coincidencias y discrepancias con la literatura previa.

El \textbf{capítulo~5} recoge, a modo de autocrítica, las principales limitaciones del trabajo
y las propuestas de mejora que de ellas se derivan. El \textbf{capítulo~6} cierra el documento
con las conclusiones —ordenadas por objetivo específico—, una reflexión sobre la relación del
TFG con la Agenda 2030 y las futuras líneas de investigación. Tras la bibliografía, los anexos
recogen el material complementario que respalda el análisis.

Como complemento a este documento, los resultados se han volcado además en un
\textbf{cuadro de mando interactivo} de elaboración propia, que permite explorar visualmente
el corpus, las empresas individuales, los temas detectados y los principales hallazgos por
pregunta de investigación.\footnote{El cuadro de mando es un documento HTML autocontenido
(\texttt{results/dashboard/index.html}) que reúne, en cinco secciones interactivas, la visión
agregada del corpus, un explorador por empresa, el mapa de temas, un comparador entre
compañías y los resultados sintetizados de las cuatro preguntas de investigación. Constituye
la salida aplicada del trabajo y se describe con mayor detalle en el capítulo de resultados.}
